\section{ISAC的历史发展}

在最早的ISAC实现方案（20世纪60年代）中，通信信息通过脉冲间隔调制（PIM）嵌入到一组雷达脉冲中，当时的雷达被用作导弹测距仪器\cite{4337601}。尽管该方案看似简单，但当时的研究人员，远在现代数字通信诞生之前，就已经意识到某些通信功能可以集成到军用雷达中。事实上，作为传感技术的重要代表，雷达的发展深受无线通信的影响，反之亦然。本节将从历史的角度概述ISAC技术的发展历程。

\subsection{雷达的早期发展}

早期的雷达由机械马达驱动，通过周期性旋转天线来搜索空间目标。然而，这类雷达面临着诸多关键挑战，例如缺乏多功能性和灵活性，以及相对容易受到干扰。鉴于此，相控阵雷达（又称电子扫描阵列雷达）应运而生，旨在克服这些缺点。
相控阵系统并非机械旋转天线，而是生成可电子控制的空间波束。二战期间，德国GEMA公司研制了第一部远程预警相控阵雷达，名为“FuMG 41/42 Mammut”（简称“Mammut”），能够探测300公里范围内8公里高度的飞行目标

“猛犸”雷达系统不仅是首个相控阵雷达系统，也是首个多天线系统，它启发了多输入多输出（MIMO）通信系统的发明。1994年，Paulraj和Kailath获得了首个MIMO通信专利，开启了3G、4G和5G无线网络的新时代。受MIMO通信技术的启发，麻省理工学院林肯实验室在十年后的2004年IEEE雷达会议上提出了共置MIMO雷达方案。在MIMO雷达中，每个天线发射独立的波形，而不是基准波形的相移版本\cite{8930009}。这使得虚拟孔径增大，与相控阵雷达相比，提高了灵活性和传感性能。自由度（DoF）和分集等概念是从 MIMO 通信理论中“借用”来的，成为了 MIMO 雷达理论基础的基石。

\subsection{雷达与通信相互启发}

雷达和通信的研究在20世纪90年代初至21世纪初开始融合。90年代，美国海军研究办公室（ONR）启动了先进多功能射频（RF）概念（AMRFC）项目，旨在通过将多个天线划分为不同的功能模块（例如雷达、通信和电子战模块）来设计集成射频前端。20世纪90年代至21世纪初兴起的集成系统与通信（ISAC）研究主要受到AMRFC及其后续项目（例如ONR资助的集成顶部系统（InTop）项目）的推动。在此期间，雷达界提出了各种ISAC方案，其总体思路是将通信信息嵌入到常用的雷达波形中。例如，的开创性工作提出了将线性调频信号与相移键控 (PSK) 调制相结合，这是第一个利用线性调频信号的 ISAC 波形设计。此后，许多研究工作开始关注利用雷达波形（例如线性调频信号和频率/相位编码波形）作为载波来调制通信数据。

正交频分复用（OFDM）是包括4G和5G在内的无线网络中的关键技术之一，在2010年代初期被发现可用于雷达传感。具体而言，在OFDM雷达中，可以通过简单的快速傅里叶变换（FFT）及其逆变换（IFFT）来有效降低随机通信数据的影响，并解耦延迟和多普勒处理。基于线性调频信号和OFDM信号的两种方案分别是“以传感为中心”和“以通信为中心”设计的典型例子。

2013年，美国国防高级研究计划局（DARPA）资助了另一个名为“雷达与通信共享频谱接入（SSPARC）”的项目，旨在从雷达频段释放部分6GHz以下频谱，供雷达和通信系统共享使用。这引出了认知无线电框架下的另一个有趣的研究课题——“雷达-通信共存（RCC）”。在RCC框架下，雷达和通信系统需要在同一频段内共存，且彼此之间不会产生过度干扰。除了RCC涉及的频谱共存和干扰管理之外，ISAC还致力于通过通用基础设施实现这两种功能的更深层次集成。

\subsection{雷达与通信的并行发展}

2010年，Marzetta在其开创性论文中提出了大规模MIMO（mMIMO）技术，该技术后来成为5G及未来网络的核心技术之一\cite{5595728}。三年后，即2013年，纽约大学无线技术中心（NYU WIRELESS）发表了关于利用毫米波（mmWave）信号进行移动通信可行性的里程碑式论文。自此，毫米波和mMIMO成为相辅相成的完美搭档。由于信号波长缩短，大规模MIMO天线阵列的物理尺寸可以大大缩小；而mMIMO阵列提供的高波束成形增益，使得毫米波信号的传输距离更远。然而，阻碍mMIMO毫米波技术大规模部署的关键挑战在于，由于需要大量的毫米波射频链路，其硬件成本和能耗都非常高。这迫使无线研究人员重新思考mMIMO系统的射频前端架构。其中，混合模拟-数字 (HAD) 结构通过精心设计的移相器（甚至开关）网络将大量天线与少量射频链路连接起来，从而成为一种可行的有前景的解决方案，进而降低成本和能耗。

在mMIMO技术诞生的同一年，文献\cite{5419124}提出了相控阵MIMO雷达的概念，旨在平衡相控阵雷达和MIMO雷达的性能。值得注意的是，MIMO雷达通过在每个天线上发射独立的波形，以牺牲阵列增益为代价，提高了自由度；相比之下，相控阵雷达通过将发射功率集中于目标方向，虽然能够获得更高的阵列增益，但自由度有所降低。正如通信领域的HAD结构一样，一个自然的思路是设计一种系统架构来弥合两者之间的差距，即通过移相器阵列将多个天线与有限数量的射频链路连接起来。这样，相控阵MIMO雷达便实现了相控阵雷达和MIMO雷达之间灵活的权衡\cite{5419124}。在极端情况下，当只有一个射频链路时，相控阵MIMO雷达就退化为相控阵雷达。另一方面，如果射频链路的数量等于天线的数量，则相位MIMO雷达等效于MIMO雷达。该文献探讨了利用mMIMO进行雷达探测的优势，其中在存在未知统计特性的干扰的情况下，仅需一次快照即可精确探测目标。

\subsection{S\&C的融合}

上述相似之处为两种技术的融合提供了明确的契机，使其能够通过单次传输实现传感和通信。事实上，雷达和通信技术已深度交织，并在过去十年中朝着同一方向发展，即高频段和大规模天线阵列，这本质上是对频谱和空间资源的更高需求。从通信角度来看，大带宽和天线阵列能够​​提升通信容量，并提供海量连接。另一方面，增加带宽和天线数量也将显著提高雷达的距离分辨率和角度分辨率，即更精确地探测更多目标或绘制复杂环境的能力。

由于雷达和通信的工作频率都达到了毫米波频段，因此它们在信道特性和信号处理方法方面也往往具有相似性。特别是，毫米波通信信道较为稀疏，且以视距（LoS）传播为主，这是因为其可用的传播路径不如6GHz以下频段丰富。因此，毫米波信道模型更符合物理几何结构，这与mMIMO技术相结合，推动了毫米波通信波束域信号处理技术的发展。这些技术包括但不限于波束训练、波束对准、波束跟踪和波束管理，所有这些都可以基于HAD结构。值得注意的是，波束域通信在一定程度上模拟了传统的雷达信号处理，其中波束训练和跟踪可以类比为目标搜索和跟踪。因此，雷达和通信之间的界限变得模糊，感知功能也不一定局限于雷达基础设施。无线基础设施和设备也可以通过无线电发射和信号传输进行感知，这构成了ISAC的技术基础和理论依据。

\begin{table}[htbp]
	\centering
	\caption{ISAC 技术演进的四个阶段}
	\label{tab:isac_history}
	\begin{tabular}{lp{5cm}p{5cm}}
		\toprule
		\textbf{演进阶段} & \textbf{核心特征} & \textbf{关键技术/手段} \\
		\midrule
		\textbf{分离 (Separation)} & 雷达与通信在频谱和硬件上完全解耦。 & 传统雷达信号处理、独立射频前端。 \\
		\addlinespace
		\textbf{共存 (Co-existence)} & 解决频谱资源匮乏问题，强调干扰抑制。 & 频谱共享技术、干扰对消算法 (RCC)。 \\
		\addlinespace
		\textbf{融合 (Integration)} & 共享发射波形与硬件资源，实现一体化设计。 & 联合波形设计、多束成形优化 (JSC/JRC)。 \\
		\addlinespace
		\textbf{网络化 (Native)} & 感知作为 6G 原生功能，实现覆盖全局的探测。 & 蜂窝感知网络 (PMN)、协同多站感知。 \\
		\bottomrule
	\end{tabular}
\end{table}
